% ============================================================
\subsection*{Modèle homogénéisé et hypothèses de réduction}

Après avoir présenté l'approche de Thiele, on revient au modèle homogénéisé
général afin d'identifier les hypothèses qui conduisent aux modèles réduits.
L'objectif n'est pas uniquement de retrouver une formulation simplifiée, mais
aussi de disposer d'un cadre permettant de lever progressivement certaines
hypothèses si les résultats obtenus avec le modèle réduit ne sont pas
cohérents avec le comportement attendu du système.

On étudie maintenant l'effet de différentes hypothèses simplificatrices sur ce système.

\subsubsection*{Simplifications du système}

On applique successivement les hypothèses suivantes :

\paragraph{1. Régime stationnaire :}
On suppose que le système a atteint un état stationnaire, 
c'est-à-dire que les concentrations ne varient plus dans le temps :

\begin{equation}
\frac{\partial \avg{C_i}^l}{\partial t} = 0
\end{equation}

\eqref{eq:sys1} se réduit alors à :

\begin{equation}
\nab\!\cdot
\left[
D_i^{\mathrm{tot}} \nab\avg{C_i}^l
+
\frac{z_i F}{RT}\varepsilon_l D_i^{\mathrm{eff,mig}}
\avg{C_i}^l \nab\avg{\varphi_l}^l
-
\varepsilon_l \avg{C_i}^l \avg{v}^l
\right]
= -a_v^{\mathrm{cat}} R_i
\end{equation}


\paragraph{2. Pas de convection :}
On suppose que la vitesse du fluide est négligeable :

\begin{equation}
\avg{v}^l = 0
\end{equation}

\eqref{eq:sys1} se réduit alors à :

\begin{equation}
\nab\!\cdot
\left[
D_i^{\mathrm{tot}} \nab\avg{C_i}^l
+
\frac{z_i F}{RT}\varepsilon_l D_i^{\mathrm{eff,mig}}
\avg{C_i}^l \nab\avg{\varphi_l}^l
\right]
= -a_v^{\mathrm{cat}} R_i
\end{equation}

En l'absence d'écoulement, la dispersion hydrodynamique disparaît et
$D_i^{\mathrm{tot}} = \varepsilon_l D_i^{\mathrm{eff}}$.

\paragraph{3. Pas de migration :}
On suppose que le terme migratoire est négligeable devant le terme
diffusif, ce qui revient à négliger l'effet du champ électrique sur
le transport des espèces. Le transport est alors gouverné uniquement
par la diffusion et la réaction :

\begin{equation}
\nab\!\left(\varepsilon_l D_i^{\text{eff}} \nab\avg{C_i}^l\right) = -a_v^{\mathrm{cat}} R_i
\label{eq:diff_reaction}
\end{equation}

\paragraph{4. Potentiel solide uniforme :}
Lorsque la chute de potentiel dans la phase solide est négligeable,
le potentiel solide peut être considéré comme uniforme :

\begin{equation}
\nab\avg{\varphi_s}^s = 0
\end{equation}

La loi d'Ohm homogénéisée est alors automatiquement satisfaite et
le potentiel solide devient un paramètre uniforme du problème.

\paragraph{5. Module de Thiele :}

On adimensionne l'équation \eqref{eq:diff_reaction} en posant 
\[
\avg{C_i}^{l,*} = \frac{\avg{C_i}^l}{C_i^{\text{ref}}},
\qquad 
\zeta = \frac{z}{l_e},
\]
où $l_e$ est la longueur caractéristique macroscopique.

En supposant une cinétique du premier ordre pour le courant interfacial,
\begin{equation}
j = k(\eta)\avg{C_i}^l ,
\end{equation}
et en utilisant la relation
\begin{equation}
R_i =
-\frac{\nu_i}{n_eF}j ,
\end{equation}
on obtient :
\begin{equation}
R_i =
-\frac{\nu_i k(\eta)}{n_eF}\avg{C_i}^l .
\end{equation}

L'équation de diffusion--réaction devient alors, en 1D :
\begin{equation}
\frac{d^2 \avg{C_i}^{l,*}}{d \zeta^2}
=
\underbrace{
-\frac{\nu_i a_v^{\mathrm{cat}} k(\eta) l_e^2}
{n_eF\,\varepsilon_l D_i^{\mathrm{eff}}}
}_{\phi^2}
\avg{C_i}^{l,*}.
\end{equation}

On retrouve donc le module de Thiele :
\begin{equation}
\boxed{
\phi^2
=
-\frac{\nu_i a_v^{\mathrm{cat}} k(\eta) l_e^2}
{n_eF\,\varepsilon_l D_i^{\mathrm{eff}}}
}
\end{equation}

\begin{itemize}
    \item $\phi \ll 1$ : la diffusion est rapide devant la réaction, 
    la concentration est uniforme dans l'électrode.
    \item $\phi \gg 1$ : la réaction est rapide devant la diffusion, 
    la concentration chute rapidement près de l'interface.
\end{itemize}

La solution analytique de cette équation ainsi que l'interprétation
physique du module de Thiele ont été présentées dans la section
consacrée à l'approche de Thiele.


\bigskip
Les différentes hypothèses introduites permettent de relier le système
homogénéisé général à des modèles réduits plus simples~\cite{dao3SimplificationOrder2012, schmidtUnderstandingDeviationsSpatially2021}. En particulier,
le modèle de Thiele apparaît comme un cas limite du formalisme obtenu
par homogénéisation lorsque seuls les phénomènes de diffusion et de
réaction sont conservés.